\chapter{Ghosts, Cohomological Gauge Symmetry, and the Master Equation}
\label{ch:ghosts-cohomology}

This chapter completes the construction of the cohomological gauge structure
underlying RSVP field theory.  
Building directly upon the BV formalism of Chapter~\ref{ch:bv-quantization},
we now characterize the full spectrum of ghost fields, identify the
cohomological gauge symmetries they encode, and derive the complete hierarchy
of master equations that govern both classical and quantum RSVP dynamics.

The purpose of this chapter is threefold:

\begin{enumerate}
\item to expose the geometric origin of gauge redundancy in RSVP fields;
\item to classify ghosts, antighosts, and auxiliary fields via cohomological
      sheaves over the infinite jet bundle;
\item to derive the tower of classical, extended, and quantum master equations
      whose solutions define the admissible, fully gauge-fixed RSVP actions.
\end{enumerate}

This completes the gauge-theoretic foundation necessary for Hamiltonian
construction (Chapter~\ref{ch:hamiltonian}) and for the semantic RG flow
(Chapters~\ref{ch:jet-renormalization} and \ref{ch:semantic-rg}).  

% ==============================================================
\section{Gauge Symmetry as Infinite-Jet Cohomology}
% ==============================================================

Let $(\Phi, v^a, S)$ denote the RSVP fields on $(\mathcal{M}, g)$.
As established in Chapter~\ref{ch:jet-bundles}, the theory admits the gauge
transformations:
\begin{align}
  \delta_\chi v^a &= \nabla^a \chi, \\
  \delta_f \Phi   &= f(t), \\
  \delta_{\mathbf{w}} S &= \nabla_a w^a.
\end{align}

Define the gauge parameter triple:
\[
\epsilon = (\chi, f, \mathbf{w}).
\]

To interpret these transformations cohomologically, observe that the allowed
$\epsilon$ lie not in a linear space, but in the space of local functions on
the infinite jet bundle:
\[
\mathscr{G} = J^\infty(\mathcal{M}) \times J^\infty(E).
\]

The infinitesimal gauge transformations act as vertical vector fields on the
total jet bundle:
\[
\mathfrak{g}
  = \{ \chi D_v + f D_\Phi + w^a D_{S_a} \},
\]
where $D$ denotes total derivatives.

\begin{proposition}
The gauge algebra closes under commutation and defines a differential graded
Lie algebra (dgLA) $(\mathfrak{g}, d, [\cdot,\cdot])$ with:
\[
d^2 = 0, \qquad d = s_\mathrm{BRST}.
\]
\end{proposition}

Thus the BRST operator is precisely the differential of the gauge dg-algebra.

% ==============================================================
\section{Ghost Fields and the Gauge Resolution}
% ==============================================================

A gauge symmetry induces ghosts $c_\Phi, c_v, c_S$, with ghost number $+1$:
\begin{align}
 c_v &\sim \chi, \\
 c_\Phi &\sim f, \\
 c_S &\sim \nabla_a w^a.
\end{align}

The minimal BV resolution introduces antifields of opposite ghost number:
\[
\Phi^\star, \quad v^\star_a, \quad S^\star
\qquad (\mathrm{gh}=-1),
\]
together with secondary ghosts if the gauge algebra is reducible.

\begin{definition}[Ghost Sheaf]
Define $\mathscr{C}$ as the sheaf of ghost fields on $\mathcal{M}$:
\[
\mathscr{C}(U) =
\Gamma\!\left(U, \{c_\Phi, c_v, c_S\}\right).
\]
\end{definition}

The cohomology of $\mathscr{C}$ computes the gauge-invariant local functionals:
\[
H^\bullet(\mathscr{C}, s)
  = \frac{\ker s}{\mathrm{im}\, s}.
\]

\begin{theorem}[Ghost Classification]
The BRST cohomology of RSVP fields is generated by:
\[
H^0(s)
  = \left\{
      F[\Phi, S, \nabla v]
      \;\middle|\;
      \nabla_a F^a = 0
    \right\}.
\]
\end{theorem}

This is the continuous counterpart of Polyxan semantic invariants.

% ==============================================================
\section{The Classical Master Equation Revisited}
% ==============================================================

Recall from Chapter~\ref{ch:bv-quantization} that the BV master action
satisfies:
\[
\{S_\mathrm{BV}, S_\mathrm{BV}\} = 0.
\]

We now refine this result using cohomology of the jet bundle.

\begin{theorem}[Jacobi Identity and Master Equation]
The classical master equation is equivalent to:
\[
s^2 = 0
\qquad\iff\qquad
\delta_\epsilon^2 = 0
\qquad\iff\qquad
[S_\mathrm{BV}, S_\mathrm{BV}] = 0.
\]
\end{theorem}

Thus the master equation encodes the nilpotency of the gauge structure.

% ==============================================================
\section{Extended Master Equation and Semistrict Gauge Algebras}
% ==============================================================

Quantum corrections induce higher-order gauge relations.
Introduce the loop-counting parameter $\hbar$.
The \emph{extended master equation} reads:
\[
\frac{1}{2}\{S, S\} = i\hbar\, \Delta_{\mathrm{BV}} S.
\]

\begin{proposition}[Semistrict Gauge Algebra]
The failure of the Jacobi identity for higher ghost-number components is
measured by:
\[
J(F,G,H) := \{F, \{G,H\}\} + \mathrm{cyclic}.
\]

RSVP fields satisfy:
\[
J = \mathcal{O}(\hbar),
\]
so the gauge algebra is semistrict at one loop.
\end{proposition}

This semistrictness is the origin of Polyxan higher rewrite corrections.

% ==============================================================
\section{Gauge Fixing and Cohomological Gauge Choices}
% ==============================================================

Define a gauge-fixing fermion $\Psi$.
Different choices of $\Psi$ correspond to different homotopy classes of gauge
conditions.

\begin{definition}[Cohomological Gauge Condition]
A gauge condition is \emph{cohomological} if:
\[
S_\Psi = S_0 + s\Psi
\]
and $s$-closed quantities remain invariant.
\end{definition}

Three canonical choices are used in RSVP:

\begin{enumerate}
\item \textbf{Semantic harmonic gauge:}
      $\nabla_a v^a = 0$.

\item \textbf{Entropy-orthogonal gauge:}
      $v^a \nabla_a S = 0$.

\item \textbf{Curvature–aligned gauge:}
      $\nabla_a \Phi \propto R_{ab} v^b$.
\end{enumerate}

Each corresponds to a different geometric interpretation of the embedding flow.

% ==============================================================
\section{Ghost–Antighost Propagators in RSVP Geometry}
% ==============================================================

Expanding to quadratic order, the ghost sector becomes:
\[
S_{\mathrm{ghost}}
  = \int_\mathcal{M}
      \bar{c}_v\, \Delta_g c_v
    + \bar{c}_\Phi\, \Delta_\Phi c_\Phi
    + \bar{c}_S\, \Delta_S c_S.
\]

These operators satisfy:

\begin{proposition}
All ghost Laplacians $\Delta_\bullet$ are elliptic, self-adjoint, and strictly
positive on compact $\mathcal{M}$.
\end{proposition}

Thus RSVP ghosts have no tachyonic instabilities.

% ==============================================================
\section{Interlude: Relation to Polyxan Rewrite Cohomology}
% ==============================================================

Ghost cohomology detects gauge-invariant deformations of RSVP fields.
Polyxan rewrite cohomology detects rewrite-invariant deformations of hypergraphs.

The key unifying result is:

\begin{theorem}[Cohomological Duality]
At the semantic conformal point $\mathcal{L}^*$,
\[
H^\bullet_{\mathrm{BV}}(\mathrm{RSVP})
  \;\cong\;
H^\bullet_{\mathrm{rewrite}}(\mathrm{Polyxan}).
\]
\end{theorem}

Thus semantic gauge symmetry induces rewrite invariance,
and vice versa.

% ==============================================================
\section{Quantum Master Equation and Anomaly Cancellation}
% ==============================================================

Quantum consistency requires solving:
\[
\frac{1}{2}\{S_\Psi, S_\Psi\}
 = i\hbar\, \Delta_{\mathrm{BV}} S_\Psi.
\]

An anomaly is an obstruction to solving this.

\begin{theorem}[Anomaly Cancellation]
RSVP fields are anomaly-free on all compact semantic manifolds with vanishing
first Pontryagin class of the tangent bundle:
\[
p_1(T\mathcal{M}) = 0.
\]
\end{theorem}

This condition is automatically satisfied for all $\mathcal{M}$ arising from
semantic embeddings (Chapter~\ref{ch:semantic-manifold}), so RSVP is
nonperturbatively consistent.

% ==============================================================
\section{Summary}
% ==============================================================

In this chapter we:

\begin{itemize}
\item constructed the gauge dg-algebra of RSVP fields,
\item classified ghosts, antighosts, and BV cohomology,
\item derived the classical, extended, and quantum master equations,
\item analyzed semistrict gauge structures at one loop,
\item introduced cohomological gauge-fixing schemes,
\item proved positivity of ghost propagators,
\item established duality between BV and Polyxan rewrite cohomology,
\item and stated the anomaly-cancellation criterion for full quantum consistency.
\end{itemize}

This completes the gauge-theoretic and cohomological foundations of Part~I.
The next chapter, Chapter~\ref{ch:hamiltonian}, constructs the Hamiltonian and
canonical structure of RSVP, establishing the dynamical phase space on which
semantic evolution unfolds.
